\documentclass[12pt]{article}
\usepackage[utf8]{inputenc}
\usepackage[T1]{fontenc}
\usepackage[margin=1in]{geometry}
\usepackage{setspace}
\usepackage{xcolor}
\usepackage[colorlinks=true, urlcolor=blue!60!black, linkcolor=blue!60!black]{hyperref}
\setlength{\parindent}{2em}
\setlength{\parskip}{1.25em}
\renewcommand{\baselinestretch}{1.0}
\begin{document}
\begin{center}
{\LARGE \textit{Cohen v. California}, 403 U.S. 15 (1971)\par}
\vspace{0.8em}
{\large Case Brief\par}
\end{center}
\vspace{1em}

\section{Memory Jogger}
A 19-year-old wore a jacket reading ``Fuck the Draft'' in a Los Angeles courthouse corridor; the Court held 5--4 that a State may not criminalize the public display of a single expletive conveying a political message, because it was not obscenity or fighting words and unwilling viewers could avert their eyes (\href{https://supreme.justia.com/cases/federal/us/403/15/}{Justia}).

\section{Detailed Case Facts}
On April 26, 1968, nineteen-year-old Paul Robert Cohen walked through a corridor of the Los Angeles County Courthouse outside Division 20 of the Municipal Court wearing a jacket bearing the plainly visible words ``Fuck the Draft.'' Women and children were present in the corridor. When Cohen entered a courtroom, he removed the jacket and stood with it folded over his arm; a police officer sent the presiding judge a note suggesting contempt, which the judge declined, and the officer arrested Cohen after he stepped back into the corridor (403 U.S. at 19 n.3).

Cohen testified that he wore the jacket to inform the public of the depth of his feelings against the Vietnam War and the draft. He engaged in no violence, made no loud or unusual noise, uttered no sound prior to arrest, and provoked no violent reaction or threat from anyone present.

\section{Procedural History}
He was convicted in Los Angeles Municipal Court of violating California Penal Code \S~415---which barred ``maliciously and willfully disturb[ing] the peace or quiet of any neighborhood or person \dots by \dots offensive conduct''---and sentenced to 30 days' imprisonment.

The California Court of Appeal affirmed (1 Cal.\ App.\ 3d 94 (1969)), construing ``offensive conduct'' as ``behavior which has a tendency to provoke others to acts of violence or to in turn disturb the peace,'' and reasoning that it was reasonably foreseeable that Cohen's jacket might cause someone to strike him or forcibly remove it. After the California Supreme Court denied review by a 4--3 vote, the U.S.\ Supreme Court noted probable jurisdiction and reversed.

\section{Judicial Votes}
The Supreme Court \textbf{reversed} the California Court of Appeal by a \textbf{5--4} vote. The case was argued on February 22, 1971, and decided on June 7, 1971 (\href{https://supreme.justia.com/cases/federal/us/403/15/}{Justia}; \href{https://www.law.cornell.edu/supremecourt/text/403/15}{Cornell LII}):

\begin{itemize}
    \item \textbf{Majority Opinion:} Justice Harlan delivered the opinion of the Court, joined by Justices Douglas, Brennan, Stewart, and Marshall.
    \item \textbf{Concurring Opinions:} None.
    \item \textbf{Dissenting Opinion:} Justice Blackmun filed a dissenting opinion, joined by Chief Justice Burger and Justice Black, and joined as to its second paragraph by Justice White.
\end{itemize}

\section{Holding}
\textbf{Questions Presented:}

May a State, consistent with the First and Fourteenth Amendments, criminally punish the peaceful public display of a four-letter expletive criticizing the military draft under a general breach-of-the-peace statute, absent incitement, obscenity, personally directed fighting words, or an intolerable invasion of substantial privacy interests?

\textbf{No. Reversed.} Absent a more particularized and compelling reason, the State may not criminalize the simple public display of a single four-letter expletive. Profanity used to convey political sentiment is protected speech---encompassing both cognitive and emotive expression---because ``one man's vulgarity is another's lyric,'' and passersby in a public courthouse corridor are not a captive audience; they can avoid offense simply by averting their eyes.

\section{Analysis of Opinions}
\subsection{Majority Opinion (Justice Harlan)}

Justice Harlan's landmark opinion proceeds in two stages: first clearing away inapplicable First Amendment exceptions, and then dismantling California's asserted authority to sanitize public vocabulary.

\begin{enumerate}
  \item \textbf{Pure Speech, Not Conduct or Location-Specific Regulation:} The conviction rested solely on the offensiveness of Cohen's words---pure ``speech,'' not separately identifiable noncommunicative conduct regulable under \textit{United States v. O'Brien}, 391 U.S. 367 (1968). Nor could California defend the conviction as a courthouse-decorum rule, because \S~415 applied statewide and gave no fair notice of location-specific restrictions.
  \item \textbf{No Categorical Exception Applies:}
  \begin{itemize}
    \item \textit{Not Obscenity:} Under \textit{Roth v. United States}, 354 U.S. 476 (1957), obscenity must be ``in some significant way, erotic.'' A crude anti-draft slogan conjures no erotic stimulation.
    \item \textit{Not Fighting Words or Hostile-Audience Incitement:} Under \textit{Chaplinsky v. New Hampshire}, 315 U.S. 568 (1942), ``fighting words'' are personally abusive epithets ``directed to the person of the hearer'' that inherently provoke violent retaliation. Cohen's slogan was directed at the Selective Service System, not any individual present, and no one was provoked to violence (see \href{https://www.law.cornell.edu/constitution-conan/amendment-1/fighting-words}{Cornell LII, Fighting Words}).
    \item \textit{No Captive Audience:} While government may shield the sanctuary of the home from intrusive speech, citizens outside the home are routinely subject to objectionable expression. Unwilling viewers in a courthouse corridor suffered no intolerable invasion of substantial privacy; they ``could effectively avoid further bombardment of their sensibilities simply by averting their eyes.''
  \end{itemize}
  \item \textbf{Rejection of State Linguistic Censorship:} Confronting California's core claim that States may excise a single scurrilous word to preserve public civility or prevent hypothetical violence, Harlan dismissed the violence rationale as ``undifferentiated fear or apprehension of disturbance'' (\textit{Tinker v. Des Moines Indep.\ Cmty.\ Sch.\ Dist.}, 393 U.S. 503, 508 (1969)). He then articulated three enduring pillars of First Amendment doctrine:
  \begin{itemize}
    \item \textit{No Principled Stopping Point:} Giving officials power to ban ``offensive'' words is inherently boundless because ``one man's vulgarity is another's lyric''; the Constitution leaves taste and style to the individual.
    \item \textit{Protection of Emotive Force:} Linguistic expression serves a dual function---cognitive \emph{and} emotive. Words are chosen as much for emotional intensity as detached logic, and the First Amendment protects both.
    \item \textit{Risk of Suppressing Ideas:} Government cannot forbid particular words without risking suppression of unpopular ideas, as officials may seize on word-censorship as a convenient pretext to silence dissent.
  \end{itemize}
\end{enumerate}

\subsection{Dissenting Opinion (Justice Blackmun)}

Blackmun offered two grounds for dissent. First (joined by Burger and Black), he dismissed Cohen's ``absurd and immature antic'' as ``mainly conduct, and little speech,'' placing it squarely within \textit{Chaplinsky}'s fighting-words sphere and calling the majority's First Amendment analysis ``misplaced and unnecessary.'' Second (joined by Burger, Black, and White), he argued that the case should be remanded to the California Court of Appeal for reconsideration in light of \textit{In re Bushman}, 1 Cal.\ 3d 767 (1970), an intervening California Supreme Court decision narrowing \S~415 to conduct creating a ``clear and present danger'' of violence. (Harlan's footnote 2 rejected remand because \textit{Bushman} had cited \textit{Cohen} approvingly.)

\subsection{Subsequent Developments}

\begin{itemize}
  \item \textbf{Trademarks and Offensive Language:} In \href{https://supreme.justia.com/cases/federal/us/588/18-302/}{\textit{Iancu v. Brunetti}}, 588 U.S. 388 (2019), the Court invalidated the Lanham Act's bar on registering ``immoral or scandalous'' trademarks (for the clothing brand ``FUCT''), reinforcing \textit{Cohen}'s core rule that government cannot penalize vulgar or offensive expression based on viewpoint or conventional taste.
  \item \textbf{Student Off-Campus Profanity:} In \href{https://supreme.justia.com/cases/federal/us/594/20-255/}{\textit{Mahanoy Area School District v. B.L.}}, 594 U.S. 180 (2021), the Court held that a public high school violated the First Amendment by suspending a cheerleader for an off-campus Snapchat post reading ``Fuck school fuck softball fuck cheer fuck everything.'' Both the majority and Justice Alito's concurrence (explicitly citing \textit{Cohen} and \textit{Brunetti}) reaffirmed that off-campus vulgar criticism receives full First Amendment protection absent substantial disruption.

\end{itemize}

\section{Examples: Future Applications}
\begin{enumerate}

\item \textbf{Same-Side (Speech Protected): A Profane Protest Sign Outside City Hall.} During a rally on the plaza outside city hall, a resident holds a sign using the same four-letter word to condemn the mayor's housing policy. Families walking to a nearby park see it, and police arrest the resident under an ordinance banning offensive language in public places. As in \textit{Cohen}, the sign targets a policy and a public official rather than any individual present, so it is not fighting words, it is not erotic and so not obscene, and passersby in a public square can look away. The State may not punish the word merely because it offends, and the arrest violates the First Amendment.

\item \textbf{Same-Side (Speech Protected): A Ban on Profane Bumper Stickers.} A State makes it a misdemeanor to display on any vehicle a sticker containing profanity, and police cite a driver whose sticker uses an expletive to protest a war. Other motorists and pedestrians who see the sticker are outside the home and can avert their eyes, so they are not a captive audience. The law reaches emotive political expression that \textit{Cohen} protects, gives officials the boundless power over vocabulary Harlan rejected because one man's vulgarity is another's lyric, and is not tied to any setting that would justify special decorum rules. The citation cannot stand.

\item \textbf{Opposite-Side (Speech Punishable): A Face-to-Face Epithet Inviting a Fight.} During a heated argument on a sidewalk, a man steps within inches of a rival, uses the same expletive together with personal slurs aimed at him, and dares him to throw a punch. He is convicted under a statute narrowly construed to reach only words likely to provoke the average listener to immediate violence. Unlike Cohen's jacket, this speech is a personally abusive epithet directed to the person of the hearer, the core of \textit{Chaplinsky}'s fighting-words category that \textit{Cohen} preserved. The conviction is consistent with the First Amendment.

\item \textbf{Opposite-Side (Restriction Upheld): Profanity Broadcast Into the Home.} A radio station airs, in the early afternoon, a comedy routine that repeats several vulgar words, and a parent driving with a young child complains. The federal regulator sanctions the station under its indecency rules. Unlike the corridor in \textit{Cohen}, broadcast speech reaches listeners in the privacy of the home and is uniquely accessible to children, so a warning or averting one's eyes cannot fully protect them. In \textit{FCC v.\ Pacifica Foundation}, 438 U.S.\ 726 (1978), the Court upheld a sanction on just these facts, showing that \textit{Cohen}'s rule depends on the medium and setting.

\item \textbf{Fence-Sitter (Close Call): A Courtroom Ban on Profane Apparel.} A city enacts a narrow ordinance prohibiting ``any person inside a municipal courtroom or jury-assembly room from wearing apparel displaying profanity.'' A protester wears a jacket reading ``Fuck the Police'' into the municipal courtroom while court is in session and refuses the bailiff's order to remove it. Distinguish \textit{Cohen} on three fronts. First, unlike California's statewide breach-of-the-peace statute, this ordinance is location-specific and gives clear notice (addressing Harlan's notice objection at 403 U.S. at 19). Second, a courtroom in session is a nonpublic forum where the government may impose reasonable, viewpoint-neutral decorum rules (and jurors/litigants compelled to attend are far closer to a true captive audience than corridor passersby). Third, however, the challenger will invoke \textit{Cohen} and \textit{Brunetti} to argue that banning ``profanity'' still invites arbitrary, viewpoint-skewed enforcement unless strictly tied to courtroom disruption. The case is close because \textit{Cohen} expressly left open location-specific decorum rules but condemned giving officials power to police offensive words.

\end{enumerate}

\section{Critique}
\textbf{Civility and the Public Square:} Communitarian critics contend that Harlan's ``avert your eyes'' standard undervalues the collective interest in maintaining basic norms of public decorum---especially in solemn civic spaces like courthouses or around children---and erodes civic discourse by treating every norm of restraint as a slippery slope to censorship.

\textbf{Line-Drawing Tension:} Critics note that the Court itself later drew content-based lines around profanity in specific settings (broadcast radio in \textit{Pacifica}; on-campus school assemblies in \textit{Bethel School District No.\ 403 v. Fraser}, 478 U.S. 675 (1986)), suggesting that \textit{Cohen}'s skepticism of principled distinctions is more contextual than absolute.

\section{Key Quotations}
\begin{quote}
``Those in the Los Angeles courthouse could effectively avoid further bombardment of their sensibilities simply by averting their eyes.''
\end{quote}
\noindent --- 403 U.S. at 21.

\begin{quote}
``For, while the particular four-letter word being litigated here is perhaps more distasteful than most others of its genre, it is nevertheless often true that one man's vulgarity is another's lyric. Indeed, we think it is largely because governmental officials cannot make principled distinctions in this area that the Constitution leaves matters of taste and style so largely to the individual.''
\end{quote}
\noindent --- 403 U.S. at 25.

\begin{quote}
``In fact, words are often chosen as much for their emotive as their cognitive force. We cannot sanction the view that the Constitution, while solicitous of the cognitive content of individual speech, has little or no regard for that emotive function which, practically speaking, may often be the more important element of the overall message sought to be communicated.''
\end{quote}
\noindent --- 403 U.S. at 26.

\begin{quote}
``Finally, and in the same vein, we cannot indulge the facile assumption that one can forbid particular words without also running a substantial risk of suppressing ideas in the process. Indeed, governments might soon seize upon the censorship of particular words as a convenient guise for banning the expression of unpopular views.''
\end{quote}
\noindent --- 403 U.S. at 26.


\end{document}
